\documentclass[10pt,a4paper]{amsart}
\usepackage[margin=19mm]{geometry}
\usepackage{amsmath,amssymb,mathtools}
\usepackage[scr=rsfs,cal=euler]{mathalfa}
\usepackage{xfrac}
\usepackage[unicode,hidelinks,bookmarks=false]{hyperref}
\newcommand{\ie}{i.e.\ }
\newcommand{\cf}{cf.\ }
\newcommand{\resp}{respectively\ }
\newcommand{\Subsection}{\subsection}
\providecommand{\nobreakdash}{\nobreak\hbox{-}}
\setlength{\emergencystretch}{2em}
\multlinegap0pt
\usepackage{bm}
\usepackage{bbm}
\usepackage{dsfont}
\usepackage{xspace}
\usepackage{cancel}
\usepackage{mathtools}
\usepackage{amsthm}
\makeatletter
\@ifundefined{assumption}{\newtheorem{assumption}{Assumption}}{}
\@ifundefined{theorem}{\newtheorem{theorem}{Theorem}}{}
\makeatother
\title[An efficient iteration method to reconstruct the drift term ]{An efficient iteration method to reconstruct the drift term from the final measurement}
\author{Dakang Cen, Wenlong Zhang,, Zhidong Zhang}
\date{Source: arXiv:2510.10940v1 (CC BY 4.0)}
\begin{document}
\maketitle

\noindent\emph{Source and licence.} An efficient iteration method to reconstruct the drift term from the final measurement. Version arXiv:2510.10940v1 (CC BY 4.0), distributed under the Creative Commons Attribution 4.0 International licence, as recorded in the arXiv licence metadata for this exact version and shown on the abstract page https://arxiv.org/abs/2510.10940v1. Full rights notice: licenses/arxiv-2510.10940-CC-BY-4.0.txt.\par\smallskip

\noindent\textit{Editorial scope.} Counted unit: the Theorem whose source label is \texttt{Th-sto-1} in the article (source0.tex lines 548--558, source label \texttt{Th-sto-1}), delivered together with the article's own complete original proof of it (source0.tex lines 559--638, one \texttt{proof} environment of 80 lines). No mathematics is written, repaired or re-derived: statement and proof are the article's own text line for line. Required context retained uncounted: ass-1 (lines 541--546). Every definition, display and label of those lines is retained unchanged. Omitted from this package and not redistributed: 1--540, 547--547, 639--921. Nothing inside a retained excerpt is omitted or altered: every retained body is delivered exactly as the article's lines are written, so no omission receipt of any kind accompanies this package. Citations: the retained excerpts carry no citation command at all, so no bibliography entry of the article is transcribed and no reference is redistributed. Reference adaptation: none was needed and none was made; every reference inside the delivered unit resolves inside it, so no unresolved reference or citation is produced. Printed slips kept literal: no printed word, symbol, hypothesis or number of the counted statement or of its proof was altered. Layout and class: the article's own class is \texttt{article} with packages including amsmath, amssymb, amsfonts, amsthm, authblk, color, graphicx, enumitem, mathrsfs, indentfirst, hyperref, cleveref, geometry, algorithm; the wrapper is the lane's pinned \texttt{amsart} editorial wrapper at 10pt on A4, which restates the theorem-like environments the retained text needs and adds the article's own macro definitions that the retained text needs (none), with extra wrapper packages (bm, bbm, dsfont, xspace, cancel, mathtools). The delivered numbering is the unit's own. Assets: no retained span contains a figure environment, a graphics-inclusion command, a PDF-page inclusion, a table or a TikZ picture; no image file is copied into this package, the package declares no asset dependency and no image file is redistributed with it. Licence: version arXiv:2510.10940v1 of this article is distributed under the Creative Commons Attribution 4.0 International licence, as recorded in the arXiv licence metadata for this exact version and shown on the abstract page https://arxiv.org/abs/2510.10940v1. A full rights notice is delivered at licenses/arxiv-2510.10940-CC-BY-4.0.txt. Characters: every retained excerpt is pure ASCII, so the delivered engine, XeLaTeX with its default Latin Modern text font, reports no missing character.
\begin{assumption}\label{ass-1} We assume that the eigenvalues, $\rho_1\geq\rho_2\geq\cdots$, of the eigenvalue problem
\begin{align*}
(Q\varphi,Qu)=\rho(\varphi,u),~~\forall u\in R^n,    
\end{align*}
satisfy that $\rho_k\leq Ck^{-\alpha}$, $k=1,2,\cdots,n$ and the corresponding eigenvectors form an orthonormal basis respect to the inner product.
\end{assumption}

\begin{theorem}\label{Th-sto-1}
Let Assumption \ref{ass-1} be fulfilled, and $\hat{p^*}$ be the unique solution of 
\begin{align}\label{Q-mini-1}
\min\limits_{\hat{g}\in R^n} \|Q\hat{g}-g^\delta\| + \lambda\|\hat{g}\|^2.
\end{align}
Then there exists constants $\lambda_0>0$ and $C>0$ such that for any $\lambda\leq\lambda_0$,
\begin{align*}
\mathbb{E}\big[\|\hat{g}^*-g^*\|^2\big]\leq C\lambda\|g^*\|^2+C\frac{\sigma^2}{\lambda^{1/\alpha}},
\end{align*}
where $q=Q\bar{p^*}+\epsilon$, $\Gamma\bar{p}=\bar{p^*}$.
\end{theorem}

\begin{proof} 
By deriving the necessary condition of the quadratic minimization (\ref{Q-mini-1}), we can readily see that the unique minimizer $\hat{p^*}$ satisfies the
variational equation
\begin{align}\label{Sto-err-1}
(Q\hat{p^*},Qu)+\lambda(\hat{p^*},u)=(q^e,Qu),~~\forall u\in R^n.
\end{align}
For any $u\in R^n$, we introduce $\| u \|_{\lambda}^2:=\|Qu\|^2+\lambda\|u\|^2$. By taking $u=\hat{p^*}-\bar{p^*}$ in (\ref{Sto-err-1}), one has
\begin{align}\label{Sto-err-3}
\| \hat{p^*}-\bar{p^*}\|_{\lambda}\leq \lambda^{1/2}\|\bar{p^*}\|+\sup\limits_{u\in R^n}\frac{|(\epsilon+e,Qu)|}{\| u \|_{\lambda}}.
\end{align}

Let $G=\sqrt{Q^TQ}$, $G$ is a positive definite matrix. Denote its eigenvalues system $(\hat{\rho_k},\varphi_k)$, where $\{\varphi_k\}_{k=1}^n$ is an orthonormal basis.
For $\forall u\in R^n$, $u=\sum_{k=1}^n(u,\varphi_k)\varphi_k$, $Qu=\sum_{k=1}^n(u,\varphi_k)Q\varphi_k$. It gives that
\begin{align}\label{Sto-err-4}
(u,u)=\sum_{k=1}^n(u,\varphi_k)^2(\varphi_k,\varphi_k).
\end{align}
\begin{align}\label{Sto-err-5}
(Qu,Qu)
&=\bigg(\sum_{k=1}^n(u,\varphi_k)Q\varphi_k,\sum_{k=1}^n(u,\varphi_k)Q\varphi_k\bigg)\\ \nonumber
&=\sum_{k=1}^n(u,\varphi_k)^2(G\varphi_k,G\varphi_k)\\ \nonumber
&=\sum_{k=1}^n(u,\varphi_k)^2\hat{\rho_k}^2(\varphi_k,\varphi_k),
\end{align}
where we use $(Q\varphi_i,Q\varphi_j)=(G\varphi_i,G\varphi_j)=0$, $i\neq j$. Denote 
$u_k=(u,\varphi_k)$, $\rho_k=\hat{\rho_k}^2$. Let $(\varphi_k,\varphi_k)=1/\rho_k$, $k=1,\cdots,n$. We get $(u,u)=\sum_{k=1}^n\rho_k^{-1}u_k^2$, $(Qu,Qu)=\sum_{k=1}^nu_k^2$.
From (\ref{Sto-err-4}) and (\ref{Sto-err-5}), $\| u\|_{\lambda}^2=\sum_{k=1}^n(\lambda\rho_k^{-1}+1)u_k^2$.

By the Cauchy-Schwarz inequality, it holds that
\begin{align*}
(\epsilon+e,Qu)^2
&=\bigg(\sum_{i=1}^n(e_i+\epsilon_i)\bigg(\sum_{k=1}^nu_k(Q\varphi_k)_i\bigg)\bigg)^2\\
&=\bigg(\sum_{k=1}^nu_k\bigg(\sum_{i=1}^n(e_i+\epsilon_i)(Q\varphi_k)_i\bigg)\bigg)^2\\
&\leq \sum_{k=1}^n(1+\lambda\rho_k^{-1})u_k^2\sum_{k=1}^n(1+\lambda\rho_k^{-1})^{-1}\bigg(\sum_{i=1}^n(e_i+\epsilon_i)(Q\varphi_k)_i\bigg)^2\\
&\leq \| u\|_{\lambda}^2\sum_{k=1}^n(1+\lambda\rho_k^{-1})^{-1}\bigg(\sum_{i=1}^n(e_i+\epsilon_i)(Q\varphi_k)_i\bigg)^2,
\end{align*}
that is
\begin{align*}
\frac{(\epsilon+e,Qu)^2}{\| u\|_{\lambda}^2}
\leq \sum_{k=1}^n(1+\lambda\rho_k^{-1})^{-1}\bigg(\sum_{i=1}^n(e_i+\epsilon_i)(Q\varphi_k)_i\bigg)^2.
\end{align*}
For $\big(\sum_{i=1}^n(e_i+\epsilon_i)(Q\varphi_k)_i\big)^2$, we have
\begin{align*}
\bigg(\sum_{i=1}^n(e_i+\epsilon_i)(Q\varphi_k)_i\bigg)^2
&=\bigg(\sum_{i=1}^ne_i(Q\varphi_k)_i+\sum_{i=1}^n\epsilon_i(Q\varphi_k)_i\bigg)^2\\
&\leq 2\bigg(\sum_{i=1}^ne_i(Q\varphi_k)_i\bigg)^2+2\bigg(\sum_{i=1}^n\epsilon_i(Q\varphi_k)_i\bigg)^2\\
&\leq 2\bigg(\sum_{i=1}^ne_i(Q\varphi_k)_i\bigg)^2+2\sum_{i=1}^n\epsilon_i^2\sum_{i=1}^n(Q\varphi_k)_i^2\\
&\leq 2\bigg(\sum_{i=1}^ne_i(Q\varphi_k)_i\bigg)^2+2\sum_{i=1}^n\epsilon_i^2,
\end{align*}
where $\sum_{i=1}^n(Q\varphi_k)_i^2=(Q\varphi_k,Q\varphi_k)=1$.
Then, we get
\begin{align*}
\mathbb{E}\bigg[\bigg(\sum_{i=1}^n(e_i+\epsilon_i)(Q\varphi_k)_i\bigg)^2\bigg]
&\leq 2\mathbb{E}\bigg[\bigg(\sum_{i=1}^ne_i(Q\varphi_k)_i\bigg)^2\bigg]+2\sum_{i=1}^n\epsilon_i^2\\&
\leq 2\sum_{i=1}^n\mathbb{E}[e_i^2](Q\varphi_k)_i^2+2\sum_{i=1}^n\epsilon_i^2\\
&\leq 2\sigma^2\sum_{i=1}^n(Q\varphi_k)_i^2+2\sum_{i=1}^n\epsilon_i^2\\
&\leq 2\sigma^2+2\sum_{i=1}^n\epsilon_i^2,
\end{align*}
the last estimate holds from the fact random variables $\{e_i\}_{i=1}^n$ are independent and identically distributed, i.e., $\mathbb{E}[e_ie_j]=0$, $i\neq j$. Hence, one has
\begin{align*}
\mathbb{E}\bigg[\frac{(\epsilon+e,Qu)^2}{\| u\|_{\lambda}^2}\bigg]
\leq 2\sum_{k=1}^n(1+\lambda\rho_k^{-1})^{-1}\bigg(\sigma^2+\sum_{i=1}^n\epsilon_i^2\bigg).
\end{align*}
By Assumption \ref{ass-1}, it is easy to verify that
\begin{align*}
\sum_{k=1}^n(1+\lambda\rho_k^{-1})^{-1}
&\leq \sum_{k=1}^n(1+\lambda k^{\alpha})^{-1}\\
&\leq \int_1^{\infty}(1+\lambda t^{\alpha})^{-1}dt\\
&\leq \lambda^{-1/\alpha}\int_{\lambda^{1/\alpha}}^{\infty}(1+ s^{\alpha})^{-1}ds\\
&\leq C\lambda^{-1/\alpha}.
\end{align*}
It arrvies 
\begin{align}\label{Sto-err-7}
\mathbb{E}\bigg[\frac{(\epsilon+e,Qu)^2}{\| u\|_{\lambda}^2}\bigg]
\leq C\lambda^{-1/\alpha}\bigg(\sigma^2+\sum_{i=1}^n\epsilon_i^2\bigg).
\end{align}
From (\ref{Sto-err-3}) and (\ref{Sto-err-7}), we have
\begin{align*}
\mathbb{E}\big[\| \hat{p^*}-\bar{p^*}\|_{\lambda}^2\big]\leq C\lambda\|\bar{p^*}\|^2+C\lambda^{-1/\alpha}\bigg(\sigma^2+\sum_{i=1}^n\epsilon_i^2\bigg).
\end{align*}
The proof is complete. 
\end{proof} 

\end{document}
